%%% PHP Tour 2014, Utiliser PostgreSQL en 2014
%%%
%%% À l'heure où le NoSQL passe de mode doucement, il est temps de se poser
%%% les bonnes questions vis à vis des technologies de bases de données à
%%% utiliser, comment et pourquoi. PostgreSQL entre de plein droit dans la
%%% case des SGBD relationnels classiques, aussi nous commencerons par
%%% étudier ce que de ces outils apportent. Puis nous ferons le tour des
%%% fonctionnalités avancées de PostgreSQL, qui le positionnent comme un
%%% élément clé de votre architecture d'application.

\documentclass{beamer}

\usepackage{tikz}
\usetikzlibrary{arrows,backgrounds,snakes,shadows}

\usepackage{minted}

\usepackage[utf8]{inputenc}

\usepackage{beamerthemesplit}
\usetheme{Boadilla}
%\setbeamertemplate{itemize items}{\checkmark}
\setbeamertemplate{itemize items}[circle]
\beamertemplatetransparentcovered

\usepackage{multicol}

\title{\textsc{Utiliser PostgreSQL en 2014}}
\subtitle{PHP Tour, 2014}
\author{\textit{Dimitri Fontaine} \texttt{dimitri@2ndQuadrant.fr}
  \linebreak
  \url{@tapoueh}}
\date{23 juin 2014}
\logo{\includegraphics[height=0.4cm]{2ndQuadrant-cross.png}}

\begin{document}

\frame{\titlepage}

\section{Introduction}

\begin{frame}
  \frametitle{Dimitri Fontaine}

  \center{\Large \textsc{Principal Consultant at 2ndQuadrant}}
  \begin{center}
    \includegraphics[height=1.1in]{2ndquadrant_logo_full_color.jpg}
  \end{center}
\end{frame}

\begin{frame}[fragile]
  \frametitle{Dimitri Fontaine}

  \begin{center}
    {\Large \textsc{PostgreSQL Major Contributor}}
  \end{center}

  \begin{columns}[c]
    \column{.5\textwidth} 

    \begin{itemize}
    \item \texttt{pgloader}
    \item \texttt{prefix}, \texttt{skytools}
    \item \texttt{apt.postgresql.org}
    \item \texttt{\textbf{CREATE EXTENSION}}
    \item \texttt{\textbf{CREATE EVENT TRIGGER}}
    \item \textit{Bi-Directional Réplication}
    \item \texttt{pginstall}
    \end{itemize}  

    \column{.5\textwidth}
    \begin{center}
      \includegraphics[height=9em]{postgres-logo.png}
    \end{center}
  \end{columns}
\end{frame}

\begin{frame}
  \frametitle{Système de Gestion de Bases de données Relationnelles}

  \begin{center}
    \includegraphics[height=2.1in]{machine-d-Anticythere-c.jpg}
  \end{center}

  \vfill
  \begin{center}
    \textit{La Fabuleuse machine d'Anticythère}
  \end{center}
\end{frame}

\begin{frame}
  \frametitle{Système de Gestion de Bases de données Relationnelles}

  \begin{center}
    \includegraphics[height=2.1in]{anticy6.jpg}
  \end{center}

  \vfill
  \begin{center}
    \textit{La Fabuleuse machine d'Anticythère}
  \end{center}
\end{frame}

\begin{frame}
  \frametitle{Système de Gestion de Bases de données Relationnelles}

  \begin{center}
    \includegraphics[height=2.1in]{IBM_Electronic_Data_Processing_Machine_-_GPN-2000-001881.jpg}
  \end{center}
\end{frame}

\begin{frame}
  \frametitle{Système de Gestion de Bases de données Relationnelles}

  \begin{center}
    \includegraphics[height=2.1in]{caution_acid.png}
  \end{center}
\end{frame}

\begin{frame}
  \frametitle{Transactions}

  \begin{center}
    \includegraphics[height=1.6in]{rollback-wordmark.png}
  \end{center}
\end{frame}

\begin{frame}[fragile]
  \frametitle{Transactions}

  \begin{minted}{postgresql}
    BEGIN;
      DROP TABLE donnees_critiques;
    ROLLBACK;
  \end{minted}
\end{frame}

\begin{frame}[fragile]
  \frametitle{Transactions}

  \begin{minted}{postgresql}
    START TRANSACTION ISOLATION LEVEL SERIALIZABLE;
      DROP TABLE donnees_critiques;
    ROLLBACK;
  \end{minted}
\end{frame}

\begin{frame}[fragile]
  \frametitle{Transactions}

  \begin{minted}{postgresql}
    BEGIN;
      CREATE INDEX ... ON ...;
      EXPAIN ANALYZE SELECT ...;
    ROLLBACK;
  \end{minted}
\end{frame}

\begin{frame}
  \frametitle{Backups cohérents à chaud}

  \begin{center}
    \includegraphics[height=2.1in]{online-backup.jpg}
  \end{center}
\end{frame}

\begin{frame}
  \frametitle{Disponibilité des services et des données}

  \begin{center}
    \includegraphics[height=2.1in]{distribution.jpg}
  \end{center}
\end{frame}

\begin{frame}[fragile]
  \frametitle{Disponibilité des services et des données}

  \begin{center}
    \begin{tikzpicture}[thick]
      \tikzstyle{node}=[minimum width=2.3cm,minimum height=1.3cm,
        rectangle,rounded corners=10pt,draw,
        fill=blue!20]

      \node[node]                 (P)  at (0,0)   {\Large Primary};
      \node[node][fill=orange!40] (A)  at (0,-3.2)  {\Large Archive};
      \node[node][fill=violet!20] (S1) at (-4,-5.2) {\Large Standby 1};
      \node[node][fill=violet!40] (S2) at (4,-5.2)  {\Large Standby 2};

      \tikzstyle{sr}=[->,dotted,very thick,>=latex]
      \tikzstyle{ar}=[->,>=latex]
      \tikzstyle{rc}=[->,dashed,>=latex]

      \draw[ar] (P) -- (A)
      node [midway,fill=white] {\texttt{archive\_command}};

      \draw[sr] (P) edge[bend right]
      node[pos=0.5,fill=white] {\textbf{Streaming Replication}}
      (S1);

      \draw[sr] (P) edge[bend left]
      node[pos=0.5,fill=white,text=violet]
      {\textbf{\textit{Streaming Replication}}}
      (S2);

      \draw[rc] (A) |- (S2)
      node [midway,fill=white,text=violet]
      {\texttt{restore\_command}};
    \end{tikzpicture}
  \end{center}
\end{frame}

\begin{frame}
  \frametitle{Cohérence des données}

  \begin{center}
    \includegraphics[height=2.1in]{bits.jpeg}
  \end{center}
\end{frame}

\begin{frame}[fragile]
  \frametitle{Type de données et contraintes}

  \begin{minted}{postgresql}
    CREATE TABLE reservation (
        during tsrange NOT NULL,
        EXCLUDE USING gist (during WITH &&)
    );
    
    CREATE TABLE circles (
      c circle,
      EXCLUDE USING gist (c WITH &&)
    );
  \end{minted}
\end{frame}

\begin{frame}
  \frametitle{La console interactive \texttt{psql}}

  \begin{center}
    \includegraphics[height=2.1in]{repl.png}
  \end{center}
\end{frame}

\begin{frame}[fragile]
  \frametitle{Types de données avancés}

\begin{columns}
\column{.7\textwidth}
\begin{minted}{postgresql}
select iprange, locid
  from geolite.blocks
 where iprange >>= '91.121.37.122';


          iprange          | locid 
---------------------------+-------
 91.121.0.0-91.121.159.255 |    75
(1 row)

Time: 1.220 ms
\end{minted}
\end{columns}
\end{frame}

\begin{frame}[fragile]
  \frametitle{IP Ranges, \texttt{ip4r}, Geolocation}

  \center{Géolocalisation et méta-données obtenues sur une jointure}
  \vfill

\begin{columns}
\column{.45\textwidth}
\begin{minted}{postgresql}
 select *
   from geolite.blocks
   join geolite.location
        using(locid)
  where iprange
            >>=
        '74.125.195.147';
\end{minted}
\column{.55\textwidth}
\begin{minted}{postgresql}
-[ RECORD 1 ]----------------------------
locid      | 2703
iprange    | 74.125.189.24-74.125.255.255
country    | US
region     | CA
city       | Mountain View
postalcode | 94043
location   | (-122.0574,37.4192)
metrocode  | 807
areacode   | 650

Time: 1.335 ms
\end{minted}
\end{columns}
\end{frame}

\begin{frame}
  \frametitle{\textsc{Structured Query Language}}

  \begin{center}
    \includegraphics[height=2.1in]{sql.png}
  \end{center}
\end{frame}

\begin{frame}[fragile]
  \frametitle{How Far is The Nearest Pub}

  \center{The \texttt{point} datatype is in-core}
  \vfill

\begin{columns}[c]
\column{.45\textwidth} 

\begin{minted}{postgresql}
   CREATE TABLE pubnames
   (
     id   bigint,
     pos  POINT,
     name text
   );
\end{minted}  

\column{.55\textwidth}
\begin{center}
  \includegraphics[height=9em]{ltlng.jpg}
\end{center}
\end{columns}
\end{frame}

\begin{frame}[fragile]
  \frametitle{Eearh Distance}

\begin{center}
  \includegraphics[height=12em]{latitude_and_longitude.jpg}
\end{center}
\end{frame}

\begin{frame}[fragile]
  \frametitle{How Far is The Nearest Pub, in Miles please.}

\begin{minted}{postgresql}
create extension cube;
create extension earthdistance;
\end{minted}  
\vfill

\begin{columns}
\column{.65\textwidth}
\begin{minted}{postgresql}
  select name,
   pos <@> point(-6.25,53.34) miles
    from pubnames
order by pos <-> point(-6.25,53.34)
   limit 3;
\end{minted}  
\column{.35\textwidth}
\begin{minted}{postgresql}
  name  | miles 
--------+-------
 Ned's  |  0.06
 Sub Lo |  0.07
 O'Neil |  0.12
(3 rows)

Time: 1.335 ms
\end{minted}  
\end{columns}
\end{frame}

\begin{frame}[fragile]
  \frametitle{Geolocation}

\begin{center}
  \includegraphics[height=12em]{geolocation.png}
\end{center}
\end{frame}

\begin{frame}[fragile]
  \frametitle{Geolocalisation et meta-données}

\begin{columns}
\column{.5\textwidth}
\begin{minted}{postgresql}
with geoloc as
 (
  select location as l
    from location
    join blocks using(locid)
   where iprange
         >>=
         '212.58.251.195'
 )
  select name,
         pos <@> l miles
    from pubnames, geoloc
order by pos <-> l
   limit 10;
\end{minted}  
\column{.5\textwidth}
\begin{minted}{postgresql}
        name        | miles 
--------------------+-------
 Blue Anchor        | 0.299
 Dukes Head         | 0.360
 Blue Ball          | 0.337
 Bell (aka The Rat) | 0.481
 on the Green       | 0.602
 Fox & Hounds       | 0.549
 Chequers           | 0.712
 Sportsman          | 1.377
 Kingswood Arms     | 1.205
 Tattenham Corner   | 2.007
(10 rows)

Time: 3.275 ms
\end{minted}  
\end{columns}
\end{frame}

\begin{frame}[fragile]
  \frametitle{Trigrams}

\begin{center}
  \includegraphics[height=12em]{trigramme.png}
\end{center}
\end{frame}

\begin{frame}[fragile]
  \frametitle{Trigrammes et autocompletion}

\begin{columns}
\column{.5\textwidth}
\begin{minted}{postgresql}
  select actor
    from products
   where actor % 'fran'
order by actor <-> 'fran'
   limit 10;
\end{minted}  
\column{.5\textwidth}
\begin{minted}{postgresql}
    actor     
--------------
 FRANK HAWKE
 FRANK BERRY
 FRANK POSEY
 FRANK HAWKE
 FRANCES DEE
 FRANK LEIGH
 FRANCES DAY
 FRANK FOSTER
 FRANK HORNE
 FRANK TOMEI
(10 rows)

Time: 2.960 ms
\end{minted}  
\end{columns}
\end{frame}

\begin{frame}[fragile]
  \frametitle{Tags et tableaux}

\begin{center}
  \includegraphics[height=12em]{wordpres-seo-categories-tags.jpg}
\end{center}
\end{frame}

\begin{frame}[fragile]
  \frametitle{PostgreSQL sait utiliser des tableaux}

\begin{columns}
\column{.7\textwidth}
\begin{minted}{postgresql}
  select tt.tid,
         array_agg(tags.rowid) tags
    from tags
         join tid_tag tt
           on tags.rowid = tt.tag
group by tt.tid
   limit 3;
\end{minted}  
\column{.3\textwidth}
\begin{minted}{postgresql}
 tid |   tags    
-----+-----------
   1 | {1,2}
   2 | {3,4}
   3 | {5,6,7,8}
(3 rows)

time: 942.074 ms
\end{minted}  
\end{columns}
\end{frame}

\begin{frame}
  \frametitle{Types de données composites}

  \begin{columns}[c]
    \column{.5\textwidth}
    \begin{center}
      \includegraphics[height=9em]{xml.png}
    \end{center}
    \column{.5\textwidth}
    \begin{center}
      \includegraphics[height=9em]{logo-json.png}
    \end{center}
  \end{columns}
\end{frame}

\begin{frame}
  \frametitle{SQL en 2014 : les \textit{window functions}}

  \begin{center}
    \includegraphics[height=2.1in]{how-to-frame-a-partition-1.jpg}
  \end{center}
\end{frame}

\begin{frame}[fragile]
  \frametitle{Finding the last counter value before \textit{reset}}

\begin{columns}
\column{.65\textwidth}
\center{\textit{Write some \alert{SQL} here}}
\column{.35\textwidth}
\begin{minted}{postgresql}
 tick | nb | max 
------+----+-----
    1 |  0 |    
    2 | 10 |    
    3 | 20 |    
    4 | 30 |    
    5 | 40 |  40
    6 |  0 |    
    7 | 20 |    
    8 | 30 |    
    9 | 60 |  60
(9 rows)
\end{minted}
\end{columns}
\end{frame}

\begin{frame}[fragile]
  \frametitle{Window Functions: \texttt{lead() over()}}

\begin{columns}
\column{.65\textwidth}
\begin{minted}{postgresql}
  select tick,
         nb,
         lead(nb) over (order by tick)
    from measures;
\end{minted}

\column{.35\textwidth}
\begin{minted}{postgresql}
 tick | nb | lead 
------+----+------
    1 |  0 |   10
    2 | 10 |   20
    3 | 20 |   30
    4 | 30 |   40
    5 | 40 |    0
    6 |  0 |   20
    7 | 20 |   30
    8 | 30 |   60
    9 | 60 |     
(9 rows)
\end{minted}
\end{columns}
\end{frame}

\begin{frame}[fragile]
  \frametitle{Window Functions and \texttt{CASE}}

\begin{columns}
\column{.65\textwidth}
\begin{minted}{postgresql}
  select tick, nb,
         case when lead(nb) over w < nb
              then nb

              when lead(nb) over w is null
              then nb

              else null
          end as max
    from measures
  window w as (order by tick);
\end{minted}
\column{.35\textwidth}
\begin{minted}{postgresql}
 tick | nb | max 
------+----+-----
    1 |  0 |    
    2 | 10 |    
    3 | 20 |    
    4 | 30 |    
    5 | 40 |  40
    6 |  0 |    
    7 | 20 |    
    8 | 30 |    
    9 | 60 |  60
(9 rows)
\end{minted}
\end{columns}
\end{frame}

\begin{frame}
  \frametitle{Un histogramme tout en SQL}

  \begin{center}
    \includegraphics[height=1.8in]{histogram.png}
  \end{center}
\end{frame}

\begin{frame}[fragile]
  \frametitle{Histogramme tout en SQL}
  
  \begin{minted}{postgresql}
with drb_stats as (
    select min(drb) as min,
           max(drb) as max
      from team_stats
),
     histogram as (
   select width_bucket(drb, min, max, 9) as bucket,
          int4range(min(drb), max(drb), '[]') as range,
          count(*) as freq
     from team_stats, drb_stats
 group by bucket
 order by bucket
)
 select bucket, range, freq,
        repeat('*', (freq::float / max(freq) over() * 30)::int) as bar
   from histogram;
  \end{minted}
\end{frame}

\begin{frame}[fragile]
  \frametitle{Histogramme tout en SQL}
  
  \begin{minted}{postgresql}
 bucket |  range  | freq  |              bar               
--------+---------+-------+--------------------------------
      1 | [10,15) |    52 | 
      2 | [15,20) |  1363 | **
      3 | [20,25) |  8832 | *************
      4 | [25,30) | 20917 | ******************************
      5 | [30,35) | 20681 | ******************************
      6 | [35,40) |  9166 | *************
      7 | [40,45) |  2093 | ***
      8 | [45,50) |   247 | 
      9 | [50,54) |    20 | 
     10 | [54,55) |     1 | 
(10 rows)

Time: 53.570 ms
\end{minted}
\end{frame}

\begin{frame}[fragile]
  \frametitle{Savoir contourner les ORMs avec les \texttt{wCTE}}
  
  \begin{minted}{postgresql}
with queue as (
   insert into queue (extension)
        select id
          from extension
         where shortname = $1
  returning id, extension
)
select q.id, e.id as ext_id,
       e.fullname, e.uri, e.description
  from      queue q
       join extension e on q.extension = e.id;    
  \end{minted}
\end{frame}

\begin{frame}
  \frametitle{Les Indexes dans PostgreSQL}

  \begin{center}
    \includegraphics[height=12em]{suffix-tree-banana.png}
  \end{center}
\end{frame}

\begin{frame}
  \frametitle{Les Indexes dans PostgreSQL}

  \begin{columns}
    \column{.5\textwidth}
    \begin{center}
      \includegraphics[height=12em]{B_tree_insertion_example.png}
    \end{center}

    \column{.5\textwidth}
    \begin{center}
      \includegraphics[height=12em]{rtree.png}
    \end{center}
  \end{columns}
\end{frame}

\begin{frame}
  \frametitle{Les Indexes dans PostgreSQL}

  \begin{center}
    \includegraphics[height=12em]{gin.png}
  \end{center}
\end{frame}

\begin{frame}
  \frametitle{Les jointures avec PostgreSQL}

  \begin{center}
    \includegraphics[height=12em]{huge-full-outer-join.jpg}
  \end{center}
\end{frame}

\begin{frame}[fragile]
  \frametitle{Les jointures avec PostgreSQL}

  \begin{minted}{postgresql}
WITH upd AS (
    UPDATE target t
       SET counter = t.counter + s.counter,
      FROM source s
     WHERE t.id = s.id
 RETURNING s.id
)
INSERT INTO target(id, counter)
     SELECT id, sum(counter)
       FROM source s LEFT JOIN upd t USING(id)
      WHERE t.id IS NULL
   GROUP BY s.id
   RETURNING t.id
  \end{minted}
\end{frame}

\begin{frame}
  \begin{center}
    \includegraphics[height=21em]{postgresql-mongodb.jpg}
  \end{center}
\end{frame}

\begin{frame}
  \begin{center}
    \textsc{\Huge PostgreSQL is YeSQL!}
    \vfill

    \includegraphics[height=9em]{postgres-logo.png}
  \end{center}
\end{frame}

\begin{frame}
  \frametitle{Questions?}

  \begin{center}
    Now is the time to ask!
    \vfill

    \includegraphics[height=9em]{fly.png}
  \end{center}
\end{frame}

\end{document}
